\chapter[Sermon 83. Of the Certainty of the Salvation of Saints, and What Blessing or Salvation Is.]{Of the Certainty of the Salvation of Saints, and What Blessing or Salvation Is.}
\label{sermon:083}
\markright{Sermon 83. Of the Certainty of the Salvation of Saints, and What ...}

And he said to me: Write, “Blessed are those who are called to the marriage supper of the Lamb.” And he said to me: “These are the true sayings of God.”

The second part of this chapter concerns the certainty of the salvation of the faithful. It explains, in the meantime, what kind of blessing the faithful receive. Enough has already been said about the marriage of the Lamb, which refers to the glory and blessing of the elect. However, many things in this life suggest doubt and seem to undermine salvation, so the wavering minds of believers are now strengthened. This doctrine is beneficial for afflicted and troubled consciences, and it refutes the teaching of sophists, who claim that humans are never assured of their salvation, because in another place the wise man says: “Man does not know whether he is worthy of love or hatred.” He spoke that on another occasion and for another purpose, as I have explained in my book on the grace of God, and so forth.

At this present, therefore, it is shown that the salvation of the faithful is most certain. For first, the angel commands the Evangelist to write. This is taken from the manner of men, who put in writing their testaments, covenants, and bargains, and then seal them, for the cause of credibility and for a perpetual memorial of the thing. And those who have such writings are of a quiet mind, and think themselves safe and assured against all deceptions and subtle practices. Therefore, to the intent that the mind of man might be quieted in the matter of salvation, he causes as it were an instrument to be written, whereby all the godly might be assured of certain salvation. The same manner of writing our Lord follows in other places in weighty matters, as we may see in the eighth and thirtieth chapters of Isaiah, and in the second chapter of Habakkuk. Wherefore it is less to be marveled why the Apostle Paul so often alleged that same out of Habakkuk: that the righteous shall live by faith. For this only testimony of God, as that which is cited out of the godly instrument, might suffice as all.

\index{Jerome, Saint}And where God openly commanded Moses and Jerome to write (concerning which we may certainly judge, and undoubtedly gather that other prophets also, apostles, and evangelists wrote not without commandment), we see the authority of the books of the Old and New Testaments among all the faithful. For they are divine, authentic; they are God’s instrument and testament, the books of God himself, which are rightly believed without any other help or confirmation. We believe the testaments and sealed writings of men; how much more ought we believe the canonical books of Scripture?

Again, it is plainly declared to John, what he should write: Blessed are those who are called to the marriage of the Lamb. Therefore, it is evident that now it is confirmed both by the divine oracle and lawful instrument, that those who are called to the Lamb’s supper will be and are blessed. This was proclaimed by divine oracle and written authoritatively. What place of doubtfulness is left? Undoubtedly, the faithful are blessed, grafted in Christ. For they are now called, to whom the gospel is preached, by which they are called to the participation of the gifts of God, but chiefly to eternal life through Christ; and those who not only hear the word of the gospel, but also receive it and believe it with their heart. For many are called, and few are chosen. For the gospel is preached to many, and the grace of God is offered in Christ, but they do not receive it. But those who, through the grace of God, do receive it with true faith, are blessed. For they are not only called to the Marriage, but also come to the marriage and enjoy that wedding supper. These things seem to be taken from the doctrine of our Savior, which he taught in Luke 14, of those who were bidden to the marriage. Read that passage.

\index{Resurrection}Nevertheless, it is declared along the way what that blessedness of the faithful is: nothing else truly than the enjoyment of the marriage supper of the Lamb. A supper is prepared when the day draws to a close. So also is full salvation given to the godly toward the end of the world, at the resurrection of the dead, as was explained in the previous sermon. And truly, all those things are altogether allegorical, representing to us a certain signification of eternal life and glory. Otherwise, we have learned from the doctrine of the Prophets and Apostles—which the ear has not heard, nor the eye has seen, nor has entered into the heart of man—that the same God himself has prepared for those who love him.

\index{Erasmus of Rotterdam}Finally, a most weighty assertion or confirmation of this is added. For he hears it uttered by an oracle from heaven. These words or sayings of God are true. They are truly true, and are from God. Or else, they are true because they are from God. Erasmus has translated: these words of God are true. And so has the vernacular translation: these sayings of God are true. By a double reason therefore are these things confirmed, which are here proposed: both because they are true, and because they are from God. Although they come to one point. For since they are from God, who is truth, they cannot but be true. Therefore let us believe these things, and leave no place to doubtfulness.

\index{John, Saint}Here is an explanation of why the scriptures and preaching delivered by men from the scriptures are not the word of God: for they are written on paper with ink, and pronounced with a human voice, and with a sound that passes away, whereas the word of God is not human, nor corruptible, nor passing away. For the heavenly oracle declares plainly that the sayings were written into the book of Saint John, and pronounced by the angel, and are true, and are God’s word. Paul also affirms this in 1 Thessalonians, chapter 2, that the word he preached about him was the very word of God. Likewise, Saint Peter in 1 Peter, chapter 1. Therefore, let inquisitive people cease raising these paradoxes and stop arguing that the word of God written and preached is not the word of God. Whatever is written or spoken that does not agree with the holy scripture of God is not truly the word of God. The minds of the faithful should rather be drawn to this point: that they believe and adhere to all words of scripture declared in their proper sense, as the most certain words of God. For otherwise, what can we trust? What after this can we have as undoubted and certain? To God be glory.
